% GENERATED FILE. Edit canonical lessons, then run npm run generate:books.
% canonical-source-hash: fnv1a64:bfe6ed848497d6c2
% canonical-lessons: TA-C555-kaittattal, TA-C555-arankam, TA-C555-virutu, TA-C555-vimarcanam, TA-C555-katapattiram

\chapter{Prizes and Reviews}
\label{ch:ta-prizes-and-reviews}
\hlchaptermodality{555}

\begin{chapteropening}
\textbf{By the end of this chapter:} \emph{I can talk about awards, reviews and audiences.}

Complete the last lesson of chapter 555: I can talk about awards, reviews and audiences.
\end{chapteropening}

\section[\texorpdfstring{\ta{கைத்தட்டல்} (kaittaṭṭal)}{கைத்தட்டல் (kaittaṭṭal)}]{\ta{கைத்தட்டல்} (kaittaṭṭal) — applause}

\label{lesson:TA-C555-kaittattal}



\begin{quote}
Before the new one: say the Tamil for to go fishing, then the Tamil for a puppet show.
\end{quote}

\subsection*{You'll want to know: \ta{கைத்தட்டல்}}
\textbf{\ta{கைத்தட்டல்}} — \emph{kaittaṭṭal} — \textquotedblleft{}applause\textquotedblright{}.

\begin{quote}
\ta{பாடல்} \ta{முடிந்தவுடன்} \ta{பெரிய} \ta{கைத்தட்டல்} \ta{கேட்டது}.
\end{quote}

\begin{quote}
\emph{pāṭal muṭintavuṭaṉ periya kaittaṭṭal kēṭṭatu.}
\end{quote}

\begin{quote}
As soon as the song ended, there was loud applause.
\end{quote}

Word for word \textquotedblleft{}hand-clapping\textquotedblright{}: \ta{கை} (\emph{kai}) + \ta{தட்டல்}.

\begin{cousinweb}
\textbf{\ta{கைத்தட்டல்}} (\emph{kaittaṭṭal}) is a native compound: \ta{கை} (\emph{kai}), \textquotedblleft{}a hand\textquotedblright{}, and \ta{தட்டல்} (\emph{taṭṭal}), \textquotedblleft{}a tapping, a clapping\textquotedblright{}, from \ta{தட்டு}, \textquotedblleft{}to tap, to pat\textquotedblright{} (DEDR 3039).
\end{cousinweb}

\begin{grammarlens}[title={What you've built}]
One more word for saying what happens if.
\end{grammarlens}

\subsection*{Your turn}
\begin{itemize}
\raggedright
  \item \emph{Say it:} \emph{kaittaṭṭal}
  \item \emph{Say it:} \emph{kaittaṭṭal}, once more
  \item \emph{Say it:} \emph{pāṭal muṭintavuṭaṉ periya kaittaṭṭal kēṭṭatu.} — As soon as the song ended, there was loud applause.
  \item \emph{Recall:} say the Tamil for to go fishing, then the Tamil for a puppet show, then say \emph{kaittaṭṭal} again
\end{itemize}

\begin{tcolorbox}[breakable,colback=teal!4,colframe=teal!35!black,title={Before you move on}]
What does \ta{கைத்தட்டல்} mean? (\textquotedblleft{}Applause\textquotedblright{}.) Say it once more.
\end{tcolorbox}
\section[\texorpdfstring{\ta{அரங்கம்} (araṅkam)}{அரங்கம் (araṅkam)}]{\ta{அரங்கம்} (araṅkam) — an auditorium; a stadium}

\label{lesson:TA-C555-arankam}



\begin{quote}
Before the new one: say the Tamil for a puppet show, then the Tamil for applause.
\end{quote}

\subsection*{You'll want to know: \ta{அரங்கம்}}
\textbf{\ta{அரங்கம்}} — \emph{araṅkam} — \textquotedblleft{}an auditorium; a stadium\textquotedblright{}.

\begin{quote}
\ta{அரங்கம்} \ta{நிரம்பினால்} \ta{வெளியே} \ta{திரை} \ta{வைப்பார்கள்}.
\end{quote}

\begin{quote}
\emph{araṅkam nirampiṉāl veḷiyē tirai vaippārkaḷ.}
\end{quote}

\begin{quote}
If the auditorium fills up, they will put a screen outside.
\end{quote}

It is the second half of \ta{திரையரங்கம்} (\emph{tiraiyaraṅkam}), a cinema.

\begin{cousinweb}
\textbf{\ta{அரங்கம்}} (\emph{araṅkam}) is borrowed from Sanskrit \emph{raṅga}, \textquotedblleft{}a place of public amusement, an arena, a stage\textquotedblright{} (CDIAL 10561); Tamil puts \ta{அ} before the initial \ta{ர}.
\end{cousinweb}

\begin{grammarlens}[title={What you've built}]
One more word for saying what happens if.
\end{grammarlens}

\subsection*{Your turn}
\begin{itemize}
\raggedright
  \item \emph{Say it:} \emph{araṅkam}
  \item \emph{Say it:} \emph{araṅkam}, once more
  \item \emph{Say it:} \emph{araṅkam nirampiṉāl veḷiyē tirai vaippārkaḷ.} — If the auditorium fills up, they will put a screen outside.
  \item \emph{Recall:} say the Tamil for a puppet show, then the Tamil for applause, then say \emph{araṅkam} again
\end{itemize}

\begin{tcolorbox}[breakable,colback=teal!4,colframe=teal!35!black,title={Before you move on}]
What does \ta{அரங்கம்} mean? (\textquotedblleft{}An auditorium; a stadium\textquotedblright{}.) Say it once more.
\end{tcolorbox}
\section[\texorpdfstring{\ta{விருது} (virutu)}{விருது (virutu)}]{\ta{விருது} (virutu) — an award}

\label{lesson:TA-C555-virutu}



\begin{quote}
Before the new one: say the Tamil for applause, then the Tamil for an auditorium.
\end{quote}

\subsection*{You'll want to know: \ta{விருது}}
\textbf{\ta{விருது}} — \emph{virutu} — \textquotedblleft{}an award\textquotedblright{}.

\begin{quote}
\ta{இந்த} \ta{நாவல்} \ta{தேர்வானால்} \ta{எழுத்தாளருக்குத்} \ta{தேசிய} \ta{விருது} \ta{கிடைக்கும்}.
\end{quote}

\begin{quote}
\emph{inta nāval tērvāṉāl eḻuttāḷarukkut tēciya virutu kiṭaikkum.}
\end{quote}

\begin{quote}
If this novel is chosen, the writer will get a national award.
\end{quote}

One receives it with \ta{பெறு}: \ta{விருது} \ta{பெறு}.

\begin{cousinweb}
\textbf{\ta{விருது}} (\emph{virutu}) is borrowed from Sanskrit \emph{biruda}, \textquotedblleft{}a title, a mark of honour\textquotedblright{} (Monier-Williams).
\end{cousinweb}

\begin{grammarlens}[title={What you've built}]
One more word for saying what happens if.
\end{grammarlens}

\subsection*{Your turn}
\begin{itemize}
\raggedright
  \item \emph{Say it:} \emph{virutu}
  \item \emph{Say it:} \emph{virutu}, once more
  \item \emph{Say it:} \emph{inta nāval tērvāṉāl eḻuttāḷarukkut tēciya virutu kiṭaikkum.} — If this novel is chosen, the writer will get a national award.
  \item \emph{Recall:} say the Tamil for applause, then the Tamil for an auditorium, then say \emph{virutu} again
\end{itemize}

\begin{tcolorbox}[breakable,colback=teal!4,colframe=teal!35!black,title={Before you move on}]
What does \ta{விருது} mean? (\textquotedblleft{}An award\textquotedblright{}.) Say it once more.
\end{tcolorbox}
\section[\texorpdfstring{\ta{விமர்சனம்} (vimarcaṉam)}{விமர்சனம் (vimarcaṉam)}]{\ta{விமர்சனம்} (vimarcaṉam) — a review; criticism}

\label{lesson:TA-C555-vimarcanam}



\begin{quote}
Before the new one: say the Tamil for an auditorium, then the Tamil for an award.
\end{quote}

\subsection*{You'll want to know: \ta{விமர்சனம்}}
\textbf{\ta{விமர்சனம்}} — \emph{vimarcaṉam} — \textquotedblleft{}a review; criticism\textquotedblright{}.

\begin{quote}
\ta{நல்ல} \ta{விமர்சனம்} \ta{வந்தால்} \ta{படம்} \ta{வெற்றி} \ta{பெறும்}.
\end{quote}

\begin{quote}
\emph{nalla vimarcaṉam vantāl paṭam veṟṟi peṟum.}
\end{quote}

\begin{quote}
If good reviews come, the film will be a success.
\end{quote}

A critic is a \ta{விமர்சகர்}. A -\ta{அம்} noun: \ta{விமர்சனத்தை}.

\begin{cousinweb}
\textbf{\ta{விமர்சனம்}} (\emph{vimarcaṉam}) is borrowed from Sanskrit \emph{vimarśana}, \textquotedblleft{}discussion, investigation\textquotedblright{} (Monier-Williams).
\end{cousinweb}

\begin{grammarlens}[title={What you've built}]
One more word for saying what happens if.
\end{grammarlens}

\subsection*{Your turn}
\begin{itemize}
\raggedright
  \item \emph{Say it:} \emph{vimarcaṉam}
  \item \emph{Say it:} \emph{vimarcaṉam}, once more
  \item \emph{Say it:} \emph{nalla vimarcaṉam vantāl paṭam veṟṟi peṟum.} — If good reviews come, the film will be a success.
  \item \emph{Recall:} say the Tamil for an auditorium, then the Tamil for an award, then say \emph{vimarcaṉam} again
\end{itemize}

\begin{tcolorbox}[breakable,colback=teal!4,colframe=teal!35!black,title={Before you move on}]
What does \ta{விமர்சனம்} mean? (\textquotedblleft{}A review; criticism\textquotedblright{}.) Say it once more.
\end{tcolorbox}
\section[\texorpdfstring{\ta{கதாபாத்திரம்} (katāpāttiram)}{கதாபாத்திரம் (katāpāttiram)}]{\ta{கதாபாத்திரம்} (katāpāttiram) — a story character}

\label{lesson:TA-C555-katapattiram}



\begin{quote}
Before the new one: say the Tamil for an award, then the Tamil for a review.
\end{quote}

\subsection*{You'll want to know: \ta{கதாபாத்திரம்}}
\textbf{\ta{கதாபாத்திரம்}} — \emph{katāpāttiram} — \textquotedblleft{}a story character\textquotedblright{}.

\begin{quote}
\ta{இந்தக்} \ta{கதையில்} \ta{எனக்குப்} \ta{பிடித்த} \ta{கதாபாத்திரம்} \ta{பாட்டிதான்}.
\end{quote}

\begin{quote}
\emph{intak kataiyil eṉakkup piṭitta katāpāttiram pāṭṭitāṉ.}
\end{quote}

\begin{quote}
My favourite character in this story is the grandmother.
\end{quote}

\ta{கதா} \textquotedblleft{}story\textquotedblright{} + \ta{பாத்திரம்} \textquotedblleft{}a vessel, a role\textquotedblright{}.

\begin{cousinweb}
\textbf{\ta{கதாபாத்திரம்}} (\emph{katāpāttiram}) is built from Sanskrit \emph{kathā}, \textquotedblleft{}a story\textquotedblright{} (CDIAL 2704), and \emph{pātra}, \textquotedblleft{}a vessel\textquotedblright{}, later \textquotedblleft{}a role\textquotedblright{} (CDIAL 8055).
\end{cousinweb}

\begin{grammarlens}[title={What you've built}]
One more word for saying what happens if.
\end{grammarlens}

\subsection*{Your turn}
\begin{itemize}
\raggedright
  \item \emph{Say it:} \emph{katāpāttiram}
  \item \emph{Say it:} \emph{katāpāttiram}, once more
  \item \emph{Say it:} \emph{intak kataiyil eṉakkup piṭitta katāpāttiram pāṭṭitāṉ.} — My favourite character in this story is the grandmother.
  \item \emph{Recall:} say the Tamil for an award, then the Tamil for a review, then say \emph{katāpāttiram} again
\end{itemize}

\begin{tcolorbox}[breakable,colback=teal!4,colframe=teal!35!black,title={Before you move on}]
What does \ta{கதாபாத்திரம்} mean? (\textquotedblleft{}A story character\textquotedblright{}.) Say it once more.
\end{tcolorbox}
